\scanpage{36}
est ramené à prouver $\unsure{\rho_u|\rho_n}\geq\rho'_u$, \inserted{puis $\deleted{\unsure{\rho}}d_{aff}\geq d'_{aff}$} Remplaçant $G$ par l'adhérence \deleted{de} $G_{1\,\mathrm{unipot}.}$ [qu'on peut supposer défini $/k(s_1)$] on est ramené \inserted{pour $\rho_u\geq\rho'_u$} à prouver que si \uline{$G_1$ est unipotent}, \uline{de $\overline{\mathrm m}$} $G_0$.
\margin{[sans restriction de lissité]}
% slanted

Facile par les points d'ordre fini.

Pour la deuxième assertion, on est ramené à prouver que \uline{si} $G_1$ \uline{est affine}, $G_0$ \uline{est affine}.
\margin{[sans restriction de lissité]}
% slanted

On est ramené au cas \deleted{\unsure{d'une}} où $S$ est \deleted{\unsure{un}} schéma \inserted{intègre} de \uline{type fini} \deleted{sur} \deleted{\unsure{le}} $\mathbb Z$, \inserted{$S$ \unsure{normal|ouvert|semi-local}, $\dim S=1$} (à préciser \ldots).

\deleted{Considérons les $ {}_nG$ et la partie \unsure{compos.} première.} On peut supposer \deleted{\unsure{le groupe}} \uline{le tore maximal de $G_1$ est splitté}. $s_0$ pt fermé, $s_1$ pt génériq. Donc en tous cas $k(s_0)$ est un corps fini.

\diagram{\begin{tikzcd}[ampersand replacement=\&,row sep=small]
G\arrow[d]\\S
\end{tikzcd}}
\figure{\begin{tikzpicture}[x=1cm,y=1cm]
\draw (0.9,1.5) .. controls (1.3,1.43) and (1.65,1.55) .. (1.85,1.51);
\draw[->] (1.76,1.49) .. controls (1.85,1.49) and (1.86,1.50) .. (1.89,1.46);
\draw[->] (1.85,1.04) -- (0.68,1.04) -- (0.68,.62);
\fill (.99,.94) circle (.035);
\draw[thick] (-.47,-1.16) .. controls (-.10,-1.04) and (.28,-.92) .. (.85,-1.10);
\fill (.13,-.99) circle (.045);
\end{tikzpicture}}

Passant à la cloture hensélienne (mais pas strictement hensélienne !) en $s_0$, d'où un $\widetilde S$, $\widetilde G$ sur $\widetilde S$, \deleted{les $ {}_n\widetilde G$} Soit $l$ distinct de la caractéristique résiduelle, $n$ une puissance variable de $l$, $ {}_n\widetilde G$ a une partie \unsure{\uline{finie}|$\widetilde C^{\rm rie}$} sur $\widetilde S$, soit $\widetilde H(n)$.

Alors $\pi_1(\widetilde S)=\pi_1(s_0)\simeq\mathbb Z$ opère sur $H(n)_{\bar s_0}=M_n$, d'où par passage à la limite des opérations de $\mathbb Z$ sur $\varprojlim M_n=M\simeq\mathbb Z_l^{2\rho_{ab}+\rho_r}$. D'une part, \deleted{ces opérations donnent} comme $M_\infty\simeq\varprojlim(\deleted{\unsure{{}_nG}}\inserted{H(n)})_{\bar s_1}$ \deleted{\unsure{ce}} [\unsure{l'immersion|immersion|opération} détermine mod une opération de $\pi_1(\widetilde S)$ sur $M_\infty$] d'où $M_\infty\subset\varprojlim_n({}_nG)_{\bar s_1}=$ \deleted{donc}
% archive sheet number: 35
